\section{Domain Signatures and Trajectories}
\label{sec:signatures}
\begin{figure*}[t]
\centering\includegraphics[width=.95\textwidth,height=.55\textheight,keepaspectratio]{figures/fig01_signatures.pdf}
\caption{\textbf{Amounts and timing progressively identify domain.} (A) Equal-model mean label prevalence, separating conversational labels \citep{kim2026society} from cognitive operations \citep{gandhi2025habits}; labels can co-occur across both families. (B) Cumulative held-out log-loss reduction from uniform eight-domain prediction (zero; 3-bit loss). Gray diamonds use only log token/segment counts. Open markers add amounts; filled markers add timing. Blue circles and ochre squares match the families in A. Whiskers are pointwise 95\% bootstrap intervals for cumulative gains, not the increments. The paired tests for each increment meet the interval and FDR criteria in all 14 model--family comparisons.}
\label{fig:signatures}
\end{figure*}
\subsection{Both amount and timing distinguish domains}
Operation amounts improve held-out domain prediction beyond length in all 14 model-by-annotation-group tests. Timing adds information beyond amount in all 14 tests (Figure~\ref{fig:signatures}). Median gains are 0.756 bits per trace for amount and 0.078 for timing. Amount gains range from 0.306 to 1.134 bits; timing gains range from 0.035 to 0.176. The timing effect is smaller, but recurs across every tested checkpoint and label group. Adding eight shuffled timing features instead slightly worsens prediction (median mean gain $-0.014$ bits); real timing exceeds both this control and within-trace position shuffling in all 14 comparisons with Bonferroni-adjusted bootstrap intervals (Appendix~\ref{app:random-timing}). Thus, domain differences concern where operations occur as well as how often they appear.

Selected contrasts make these signatures interpretable. Idea generation exceeds math in perspective shifts by 16.9 percentage points (pp) and backtracking by 9.1 pp, with both directions consistent across all seven models. Planning exceeds moral reasoning in verification by 18.7 pp; math exceeds planning in subgoals by 9.0 pp, again in all seven models. These are descriptive contrasts chosen to explain the profiles, not additional multiplicity-adjusted discoveries.

GPQA and cybersecurity are closest in the eight-rate space (2.21 pp RMS difference [2.00, 2.44]). Moral and safety are mutual nearest neighbors; math is closest to code. Idea generation is comparatively distant, with unstable nearest-neighbor identity under resampling. Shared formats can contribute to these relationships; proximity does not establish semantic equivalence (Appendix~\ref{app:similarity}).

\begin{figure*}[t]
\centering\includegraphics[width=.90\textwidth,height=.44\textheight,keepaspectratio]{figures/fig02_trajectories.pdf}
\caption{\textbf{Temporal allocation differs across domains.} Thick curves are equal-model means of per-trace prevalence in relative-position bins; thin curves show all seven models and are not confidence intervals. Math and planning increasingly emphasize verification. Idea generation emphasizes alternatives and mid-trace revision. Each operation is plotted independently; curves do not imply exclusive states or observed transitions.}
\label{fig:trajectories}
\end{figure*}
\subsection{Trajectories and illustrations across tasks}
In math, verification rises from 9.9\% of segments in the early third to 27.4\% in the late third, while subgoals decline from 17.4\% to 9.6\%. The verification increase occurs in all seven models. Planning similarly increases verification (16.1\% to 30.3\%), but starts from a much lower subgoal rate. Idea generation instead has high early and middle perspective-shift prevalence (33.7\% and 35.0\%), falling to 19.9\% late; backtracking peaks in the middle third at 14.0\% (Figure~\ref{fig:trajectories}). Individual curves reveal heterogeneity that the mean alone would hide.

Outcome-blind Qwen3.5-9B examples illustrate checking a mathematical parameterization or action preconditions, using subgoals to organize moral advice, and rejecting unoriginal research ideas. Selection uses proximity to domain-median profiles within the central half of token lengths. These examples illustrate label interpretation rather than independently validating it (Appendix~\ref{app:examples}).

\begin{figure}[t]
\centering\includegraphics[width=\columnwidth]{figures/fig03_task_demand.pdf}
\caption{\textbf{Matched changes in demand and format.} (A) Profile-distance contrasts beyond their repeated-sample noise references. (B) Raw paired demand-minus-format contrast, not the subtraction of the noise-adjusted quantities in A. 95\% item-bootstrap intervals use 160 problems per model, except 159 for DeepSeek-R1-Distill-Qwen-7B. Demand responsiveness varies substantially across checkpoints.}
\label{fig:demand}
\end{figure}
\subsection{Task demand changes profiles, unevenly across models}
Changing instructions while holding problems fixed alters the profiles beyond repeat-sampling noise for four of five models. The noise-adjusted demand distance is 0.429 [0.371, 0.488] for Qwen3.5-9B and 0.774 [0.656, 0.879] for Gemma4-12B-it, compared with 0.009 [$-0.076$, 0.095] for DeepSeek-R1-Distill-Llama-8B. Demand also exceeds format in the raw paired contrast for Qwen3.5-9B, Gemma4-12B-it, and OpenThinker3-7B. Neither DeepSeek distill has an interval excluding zero for that contrast (Figure~\ref{fig:demand}). This supports responsiveness to instructions in some checkpoints, rather than a purely topical account of every domain difference. It does not isolate the effects of training recipes.

\paragraph{From description to prediction.}
The profiles distinguish task conditions; do they also identify successful reasoning within a condition? In 94 eligible model--benchmark--annotation-family tests, neither operation amounts nor their timing provide statistically supported incremental outcome prediction after multiple-comparison correction. Models condition on length and available prompt strata; binary outcomes retain their original scale, as do continuous rubric scores. Appendix~\ref{app:outcome-spec} specifies the 112 candidate tests, 18 exclusions, effect estimates, and uncertainty. Only four of the fourteen candidate math cells are eligible; the other ten lack sufficient observations in both correctness classes. This result therefore does not establish absence of outcome information, or comprehensive nonprediction in math. Direct intervention tests address a different question from task separability.

\section{Changing Reasoning Profiles}
\label{sec:interventions}
\subsection{Interventions and evaluation}
All main intervention results use Qwen3.5-9B. A rank-eight ridge readout at decoder-layer output 16 maps standardized activations to 16 targets: eight future-operation rates and eight rates over the next four segments, within an 8,192-token window. Its regularized inverse $G$ maps target changes to residual directions. A source-domain vector is
\begin{equation}
v_s=6.75\,\frac{G(\mu_s-\bar\mu)}{\|G(\mu_s-\bar\mu)\|_2},
\label{eq:transfer}
\end{equation}
where $\mu_s$ is the domain's mean target profile and $\bar\mu$ equally averages math, GPQA, planning, idea, and moral profiles. The vector is added at all prompt and generated positions. This \emph{amplifies a domain-associated contrast}: even with an ideal linear map, a target near $\mu_t$ would move to $\mu_t+c_s(\mu_s-\bar\mu)$, not necessarily toward $\mu_s$. The readout predicts future and next-four rates; this constant intervention does not directly manipulate the relative-position trajectories in Figure~\ref{fig:trajectories}. Prompt-disjoint readout validation gives a mean 7.6\% MSE reduction over training-domain means; held-domain prediction is weaker (Appendix~\ref{app:audit}).

Objective evaluations use 199 math or 200 GPQA/planning tasks and two decoding seeds, paired by task and seed. We report accuracy at 12,288 and 32,768 tokens, alongside answer availability. Intervals bootstrap tasks after averaging seeds. Natural-completion outcomes use corrected extraction grades for completed traces and direct grades for truncated prefixes. Intervention behavior shifts use paired seed-1 annotations in the first 8,192 reasoning tokens. Details and additional estimates appear in Appendix~\ref{app:methods}.

\begin{figure*}[t]
\centering\includegraphics[width=.95\textwidth]{figures/fig04_cross_domain_transfer.pdf}
\caption{\textbf{Accuracy changes accompany changes in answer availability.} Domain-associated contrasts applied to Qwen3.5-9B at two token cutoffs. Cells show changes in correct answers (A) and answer availability (B), in percentage points; absolute unsteered percentages appear beneath each target. Outlines mark same-domain cells. All vectors use Equation~\ref{eq:transfer} and are distinct from A2. The matrix is exploratory; cells share baselines. Table~\ref{tab:transfer} reports selected paired intervals.}
\label{fig:transfer}
\end{figure*}
\begin{figure*}[t]
\centering\includegraphics[width=.95\textwidth]{figures/figA4_intervention_fidelity.pdf}
\caption{\textbf{Induced changes vary in fidelity and dose response.} (A) Cosine between the intended 16-target contrast and the achieved paired mean shift at the reference dose. Rows are source contrasts; columns are targets. (B--C) GPQA and idea contrasts applied to math, using the same 100 tasks and seed one at every dose. Solid lines show accuracy change; dotted lines show answer-availability change at 32,768 tokens. The hard-loop flag detects literal tail repetition; the unsteered rate is 11\%. These are descriptive operating points, with paired uncertainty in companion tables.}
\label{fig:fidelity}
\end{figure*}
\subsection{Fidelity varies across source--target pairs}
The moral-associated vector yields strongly aligned mean shifts on math (cosine 0.940 [0.897, 0.952]) and GPQA (0.907 [0.855, 0.935]). In contrast, GPQA-to-GPQA is poorly aligned ($-0.052$, 95\% interval [$-0.224$, $0.144$]). Math-to-GPQA has intermediate alignment (0.574 [0.500, 0.635]), whereas math-to-moral is near zero (0.033 [$-0.094$, $0.162$]; Figure~\ref{fig:fidelity}A). Equal residual norms therefore produce neither equal nor uniformly faithful behavioral changes.

These manipulation checks concern directional alignment, not exact attainment of a source profile. They expose a limitation of the evaluated inverse-readout controller; they do not establish that the representation itself cannot support useful control. In particular, the 31.2 pp short-budget loss for GPQA-to-GPQA cannot show that naturally occurring GPQA reasoning is undesirable when the intended change is not reliably induced. Appendix~\ref{app:fidelity} reports shift magnitudes and component estimates.

\subsection{Answer production and repetition change with dose}
At 12,288 tokens, the math-associated vector improves GPQA accuracy by \MathGPQAShort{} and planning by \MathPlanningShort{}. Answer availability increases by 52.5 and 20.8 pp, respectively (Figure~\ref{fig:transfer}). The moral-associated vector improves math by 12.1 pp, whereas the GPQA and idea vectors reduce math accuracy by 45.2 and 47.0 pp. The source rows are largely helpful or harmful across targets; this pattern alone provides limited evidence for task-specific matching.

At 32,768 tokens, math-to-GPQA is $-4.0$ pp and math-to-planning is $+3.0$ pp. Conversely, idea-to-math loses 77.4 pp in accuracy and 78.4 pp in answer availability. On the identical 100-task, seed-1 dose subset, the idea contrast reduces accuracy by 23, 77, and 79 pp at $0.5\times$, $1\times$, and $1.5\times$ dose (Figure~\ref{fig:fidelity}B). Its hard-loop rates are 35\%, 94\%, and 96\%, compared with 11\% unsteered; the GPQA contrast gives 27\%, 73\%, and 78\% (panel C).

The hard-loop flag is an exact textual criterion: for a response of at least 1,000 characters, its final 200 characters must occur verbatim in the preceding text. It measures literal repetition, not a latent dynamical mechanism. Together, the dose and budget comparisons show that large matrix effects involve answer production and disrupted generation. They do not yet resolve performance under a common answer-elicitation protocol. The protocol and its pending status appear in Appendix~\ref{app:audit}.

\subsection{Math-derived steering lowers paired judged originality}
On idea generation, the math-profile vector lowers pairwise originality preference under both Gemma4-31B-it and Qwen3.5-27B judges (net preferences $-0.544$ and $-0.554$). A math-style instruction instead raises originality preference ($+0.457$ and $+0.470$), while lowering feasibility preference under both judges. The two interventions cannot be treated as equivalent implementations of mathematical thinking. Separate checking/exploration experiments are secondary comparisons (Appendix~\ref{app:open}).

The math vector produces answers on all 400 allocated keyword--seed attempts, compared with 272/400 unsteered. Quality is conditional on producing an answer: the math-vector and math-prompt pairwise comparisons cover 170 and 158 of 200 keywords, respectively, with eligibility varying by decoding seed. In the separate rubric evaluation, assigning zero to unanswered attempts raises aggregate originality for the math vector even though originality decreases among paired answers. Both endpoints are relevant: an intervention may produce more answers while making those answers less original. Appendix~\ref{app:open} reports coverage and distinguishes pairwise preference from rubric scores.
